\documentclass[12pt]{article}

\usepackage{sbc-template}
\usepackage{graphicx,url}
\usepackage[utf8]{inputenc}
\usepackage[brazil]{babel}
\usepackage[latin1]{inputenc}  

     
\sloppy

\title{Controle de Concorrência em Bancos de Dados Relacionais:\\ Anomalias, Mecanismos e o Custo do Isolamento}

\author{Mateus de Souza Fonseca\inst{1}, Nome\inst{1}, Nome\inst{1}, Nome\inst{1}}


\address{Engenharia da Computação -- Centro Federal de Ensino Tecnológico Celso Suckow da Fonseca
  (CEFET)\\
Petrópolis -- RJ -- Brasil
\email{\{mateus.fonseca,emailbarauna,emailmarco,emailfernando\}@[aluno.cefet-rj.br]}
}
\begin{document} 

\maketitle

\begin{abstract}
  This paper presents the concurrency control mechanisms used by relational
  database management systems to preserve transaction isolation. The classical
  anomalies described by the SQL standard are experimentally reproduced in
  PostgreSQL, and the performance cost of each isolation level is measured.
  Results show the trade-off between correctness guarantees and throughput,
  including the write skew anomaly, which is not covered by the standard
  isolation levels.
\end{abstract}
     
\begin{resumo} 
  Este artigo apresenta os mecanismos de controle de concorrência utilizados
  por sistemas gerenciadores de bancos de dados relacionais para preservar o
  isolamento entre transações. As anomalias clássicas descritas pelo padrão SQL
  são reproduzidas experimentalmente no PostgreSQL, e o custo de desempenho de
  cada nível de isolamento é medido. Os resultados evidenciam o compromisso
  entre garantias de correção e vazão, incluindo a anomalia write skew, não
  coberta pelos níveis de isolamento do padrão.
\end{resumo}

\section{Introdução}
 
Sistemas gerenciadores de bancos de dados relacionais organizam o acesso aos
dados em torno do conceito de transação, definida como uma sequência de
operações tratadas como unidade indivisível de trabalho. 
Uma transação deve satisfazer quatro propriedades,
conhecidas pela sigla ACID: atomicidade, consistência, isolamento e
durabilidade. Enquanto atomicidade e durabilidade são asseguradas pelos
mecanismos de recuperação, a propriedade de isolamento depende de um conjunto
distinto de técnicas, denominado controle de concorrência.
 
Ambientes reais de produção raramente executam uma transação por vez. Sistemas
bancários, plataformas de comércio eletrônico e aplicações corporativas atendem
simultaneamente a dezenas ou milhares de usuários, cujas operações competem
pelos mesmos registros. Executar essas transações em série eliminaria qualquer
interferência, mas desperdiçaria a capacidade de processamento disponível e
tornaria o tempo de resposta proibitivo. O intercalamento das operações é,
portanto, uma necessidade de desempenho, não uma escolha de projeto.
 
Esse intercalamento, entretanto, introduz a possibilidade de resultados
incorretos. Quando duas transações leem o mesmo valor e cada uma calcula um novo
estado a partir dele, a segunda gravação sobrescreve o efeito da primeira, e uma
das operações é silenciosamente perdida. Anomalias como essa, descritas na
literatura como atualização perdida, leitura suja, leitura não repetível e
leitura fantasma, produzem estados de banco que nenhuma execução
sequencial das mesmas transações poderia gerar. O critério formal que
separa execuções corretas de incorretas é a serialidade. 
Garanti-la, contudo, impõe custos: as técnicas disponíveis
cobram o preço em espera, quando bloqueiam transações, ou em
retrabalho, quando as abortam e exigem reexecução.
 
Este trabalho apresenta os principais mecanismos de controle de concorrência
empregados por bancos de dados relacionais, reproduz experimentalmente as
anomalias previstas pelo padrão SQL e compara o custo de desempenho entre os
níveis de isolamento oferecidos pelo PostgreSQL. Os experimentos são conduzidos
em ambiente conteinerizado e disponibilizados publicamente, de modo que os
resultados possam ser reproduzidos por terceiros.

\bibliographystyle{sbc}
\bibliography{sbc-template}

\end{document}
